% ==============================================================================
%  Appendix B -- Equivariance proofs
%  Purpose: the step-by-step algebra that Ch. 5 states and defers. Everything
%  here has a corresponding unit test; cite the test next to each proof.
% ==============================================================================

\section{Gauge covariance of the shortest-path-averaged transport}
\label{app:proof-transport}

\TODO{Induction over $|\Delta x|_1$ in the DP recursion: each extension
step multiplies by $U_\mu(x)$ (forward) or $U^{\dg}_\mu(x - \hat e_\mu)$
(backward), both of which transform covariantly, and covariance survives
the average. Also prove the octant identity
$T_{-\Delta x}(x) = T^{\dg}_{\Delta x}(x - \Delta x)$.
Tests: \texttt{tests/test\_transport.py}.}

\section{Equivariance of the GEMHSA layer}
\label{app:proof-gemhsa}

\TODO{Layer-by-layer: channel projections commute with the adjoint action
(no color mixing); transported keys/values transform in the adjoint at
the query site; the score $\Re\Tr[Q^{\dg}\tilde K]$ is invariant (cyclic
trace + unitarity), and positional encodings act only on the invariant
score; the bilinear value $Q_v^{\dg} \tilde V$ transforms in the adjoint;
trace gate and residual preserve it. Tests:
\texttt{tests/test\_blocks.py} ($\SU{2}$ complex128, $\mathbb{Z}_2$
float64 bit-exact).}

\section{Invariance of the readout and the loss}
\label{app:proof-readout}

\TODO{Trace $\to$ MLP $\to$ reduction is invariant given adjoint inputs;
hence $\Obar(t)$, $C(\Delta)$, and the Rayleigh loss are gauge invariant.
Corollary worth stating: the attention map $\alpha(x, \Delta x)$ is gauge
invariant -- the license for Ch.~8.}
